% ============================================================================
% Fichier  : front/02_mode_emploi.tex (Mode d'emploi)
% Version  : 1.0 -- novembre 2026
% ============================================================================
% 30 h de cours magistral = 15 séances de 2 h ; TD (fiches, exercices) et TP
% (douze labs) faits par l'étudiant chez lui, hors des 30 h.
% ============================================================================

\chapter*{Mode d'emploi du cours}
\addcontentsline{toc}{chapter}{Mode d'emploi du cours}
\markboth{Mode d'emploi}{Mode d'emploi}

\section*{Structure générale}

Le cours compte un prologue et sept parties, soit dix-neuf chapitres, pour
30~heures de cours magistral réparties en quinze séances de deux heures.
Les travaux dirigés et les travaux pratiques se font à la maison, hors de
ces 30~heures~: pour les travaux dirigés, la fiche de questions de chaque
chapitre et ses exercices de fin de chapitre (environ dix heures au
total)~; pour les travaux pratiques, les douze laboratoires de NDAP SUKULU,
d'une heure chacun. Les corrigés des exercices sont réunis dans un livret
séparé.

\begin{tableaucours}{cL{6.4cm}cL{3.2cm}}{%
    \tableauHeader{Partie} &
    \tableauHeader{Thème} &
    \tableauHeader{Chap.} &
    \tableauHeader{Niveau}
}
0 & Prologue~: le 24 juin 2026 & P & introductif \\
I & Fondements~: la donnée, la personne, le traitement & 1--2 & fondamental \\
II & Le contexte international & 3--5 & fondamental \\
III & Le contexte national & 6--7 & fondamental \\
IV & Méthode~: auditer la protection des données & 8--12 & fondamental à
intermédiaire \\
V & Pratique~: de l'écart à la conformité & 13--16 & intermédiaire \\
VI & Défis avancés & 17--18 & intermédiaire à expert \\
VII & Recherche~: mesurer la conformité, articuler les normes & 19 & recherche \\
\end{tableaucours}

\section*{Les quinze séances}

\begin{tableaucours}{cL{6.8cm}L{4.2cm}}{%
    \tableauHeader{Séance} &
    \tableauHeader{Contenu} &
    \tableauHeader{Travail à la maison}
}
1 & Prologue et Chap.~1~: données, traitements, acteurs & Fiches P et 1 \\
2 & Chap.~2~: les principes du traitement & Fiche 2, lab 01 \\
3 & Chap.~3 et~4~: sources internationales et africaines, RGPD & Fiches 3
et 4 \\
4 & Chap.~5~: NIST et normes ISO & Fiche 5, lab 12 \\
5 & Chap.~6 et~7~: loi n°~2024/017, loi n°~2010/012, choix des critères &
Fiches 6 et 7, lab 11 \\
6 & Chap.~8~: conduire un audit de protection des données & Fiche 8 \\
7 & Chap.~9~: cartographier les traitements & Fiche 9, lab 02 \\
8 & Chap.~10~: évaluer le risque pour les personnes & Fiche 10, lab 03 \\
9 & Chap.~11~: auditer les droits des personnes & Fiche 11, lab 04 \\
10 & Chap.~12~: auditer la sécurité des données & Fiche 12, lab 05 \\
11 & Chap.~13 et~14~: sous-traitants, transferts internationaux & Fiches 13
et 14, labs 06 et 07 \\
12 & Chap.~15~: violations de données & Fiche 15, lab 08 \\
13 & Chap.~16~: programme de conformité et rapport & Fiche 16, lab 09 \\
14 & Chap.~17~: profilage, décisions automatisées, IA & Fiche 17, lab 10 \\
15 & Chap.~18 et~19~: secteurs sensibles, mesurer et articuler & Fiches 18
et 19 \\
\end{tableaucours}

Cette répartition est indicative~; l'enseignant l'adapte au calendrier et au
rythme du groupe.

\section*{Le gabarit des chapitres}

Chaque chapitre suit le même plan~: une épigraphe, le niveau et les
prérequis, les objectifs, une situation d'ouverture tirée de l'un des trois
cas, des sections numérotées, ce qu'il faut retenir, des points de
contrôle, trois à quatre exercices gradués, des lectures, et une
transition.

\section*{Les encadrés}

\begin{tableaucours}{L{3.6cm}L{8.6cm}}{%
    \tableauHeader{Encadré} &
    \tableauHeader{Rôle}
}
Objectifs & Ce que l'on doit savoir faire à la fin du chapitre \\
Définition & Les notions, en gras dans le texte et indexées \\
Violation typique & Une erreur fréquente, et la règle qui l'évite \\
Information & Un complément, un renvoi, une nuance \\
Mission (NYANGA, MBOLO, NOUN PAY) & Le cas~: les constatations, puis
l'analyse, avec un niveau de confiance \\
Traitement & La fiche d'un traitement audité~: finalité, base légale,
données, destinataires, durée, ce que l'audit en a établi \\
Démarche d'audit & Les étapes d'une procédure d'audit, dans l'ordre \\
Lecture & Un décompte de cotes \OUI{}~/ \NON{}~/ \NEK{} et sa lecture \\
Métier & Ce que la notion change pour chacun des métiers du cours \\
Texte juridique & Un article cité dans son texte officiel, et sa portée \\
Ce qu'il faut retenir & Les faits et les règles du chapitre \\
Exercice & Trois niveaux de difficulté~: \niveaubase, \niveaumoyen,
\niveauexpert \\
\end{tableaucours}

\section*{Les conventions}

\begin{itemize}
    \item Les cotes~: \OUI{} (conformité établie par une preuve), \NON{}
          (écart établi par une preuve), \NEK{} (connaissance insuffisante
          pour trancher). Chaque exigence auditée reçoit l'une de ces trois
          cotes.
    \item Constatation et analyse sont toujours séparées~; l'analyse porte
          un niveau de confiance (faible, moyenne, élevée).
    \item « La loi » désigne la loi n°~2024/017 du 23~décembre~2024
          relative à la protection des données à caractère personnel au
          Cameroun~; ses articles sont cités sans autre précision.
    \item Les textes~: seuls les articles lus dans leur version officielle
          sont cités entre guillemets~; les normes sous licence sont
          paraphrasées~; le degré de vérification figure à
          l'annexe~\ref{ann:textes}.
    \item Les cas sont fictifs~: noms en majuscules, montants en francs
          CFA, villes réelles. Les données personnelles des exemples sont
          inventées.
\end{itemize}
